\documentclass[11pt,a4paper]{article}
\usepackage[utf8]{inputenc}
\usepackage[vietnamese]{babel}
\usepackage{amsmath,amssymb}
\usepackage{geometry}
\geometry{a4paper, total={180mm,260mm}, left=15mm, right=15mm, top=18mm, bottom=18mm}
\usepackage{tikz}
\usetikzlibrary{shapes.geometric, arrows.meta, positioning, calc}
\usepackage{tcolorbox}
\tcbuselibrary{skins,breakable}
\usepackage{xcolor}
\usepackage{multicol}
\usepackage{array}
\usepackage{fontawesome}
\usepackage{fancyhdr}

\pagestyle{fancy}
\fancyhf{}
\renewcommand{\headrulewidth}{0.5pt}
\renewcommand{\footrulewidth}{0.5pt}
\fancyhead[L]{\small\bfseries\color{teal}\faBook\ SỔ TAY KIẾN THỨC TOÁN 3}
\fancyhead[R]{\small\bfseries\color{teal}Lớp Toán Thầy Lĩnh -- loptoanthaylinh.vn}
\fancyfoot[C]{\small\color{gray}-- Sổ tay kiến thức: Bài 70. Nhân số có năm chữ số với số có một chữ số --}

\definecolor{primary}{HTML}{005F73}
\definecolor{secondary}{HTML}{0A9396}
\definecolor{accent}{HTML}{EE9B00}
\definecolor{danger}{HTML}{AE2012}
\definecolor{boxbg}{HTML}{F4F6F8}

\begin{document}

% --- BANNER TIÊU ĐỀ ---
\begin{tcolorbox}[
  enhanced,
  colback=white,
  colframe=primary,
  arc=3mm,
  boxrule=1.2pt,
  left=4mm, right=4mm, top=2.5mm, bottom=2.5mm,
  drop shadow=black!10!white
]
  \centering
  {\Large\bfseries\color{primary}\faBookmark\ BÀI 70: NHÂN SỐ CÓ 5 CHỮ SỐ VỚI SỐ CÓ 1 CHỮ SỐ}\\[1.5mm]
  {\small\itshape\color{black!75}Kĩ thuật đặt tính, thứ tự nhân từ phải sang trái, nhân có nhớ không quá 2 lần và tính nhẩm tròn chục nghìn trên 1 trang A4}
\end{tcolorbox}

\vspace{-2mm}


\noindent
\begin{minipage}[t]{0.49\textwidth}\vspace{0pt}
\begin{tcolorbox}[
  enhanced,
  colback=blue!4!white,
  colframe=blue!75!black,
  arc=2mm, boxrule=0.8pt,
  left=2mm, right=2mm, top=1.8mm, bottom=1.8mm,
  title={\footnotesize\bfseries\color{white}\faTimes\ 1. Quy trình Đặt tính rồi tính}
]
{\footnotesize
\begin{itemize}\setlength{\itemsep}{1.5pt}\setlength{\leftskip}{-2.5mm}
  \item \textbf{Đặt tính:} Viết thừa số có 1 chữ số dưới thừa số có 5 chữ số, thẳng cột hàng đơn vị. Viết dấu $\times$, kẻ vạch ngang.
  \item \textbf{Thực hiện phép nhân:} Lấy thừa số thứ hai nhân lần lượt với từng chữ số của thừa số thứ nhất từ \textbf{phải sang trái}:
  $$\text{ĐV} \to \text{Chục} \to \text{Trăm} \to \text{Nghìn} \to \text{Chục nghìn}$$
  \item \textbf{Quy tắc nhớ:} Nhân trước, rồi cộng số nhớ (nếu có). Viết chữ số đơn vị của kết quả, nhớ chữ số chục sang hàng tiếp theo.
\end{itemize}
}
\end{tcolorbox}
\end{minipage}
\hfill
\begin{minipage}[t]{0.49\textwidth}\vspace{0pt}
\begin{tcolorbox}[
  enhanced,
  colback=teal!4!white,
  colframe=teal!75!black,
  arc=2mm, boxrule=0.8pt,
  left=2mm, right=2mm, top=1.8mm, bottom=1.8mm,
  title={\footnotesize\bfseries\color{white}\faKey\ 2. Nhân nhẩm số tròn chục nghìn}
]
{\footnotesize
\begin{itemize}\setlength{\itemsep}{1.5pt}\setlength{\leftskip}{-2.5mm}
  \item \textbf{Quy tắc nhẩm:}
  Lấy chữ số hàng chục nghìn nhân với số có một chữ số, rồi viết thêm bốn chữ số $0$ vào tận cùng bên phải tích.
  \item \textbf{Ví dụ nhẩm:}
  $$20\,000 \times 4 = \mathbf{80\,000}$$
  (Nhẩm: $2\text{ chục nghìn} \times 4 = 8\text{ chục nghìn}$).
  $$50\,000 \times 2 = \mathbf{100\,000}$$
\end{itemize}
}
\end{tcolorbox}
\end{minipage}

\vspace{-2.5mm}
% --- KHỐI 3: SƠ ĐỒ TƯ DUY TIKZ ---
\noindent
\begin{tcolorbox}[
  enhanced,
  colback=orange!4!white,
  colframe=orange!80!black,
  arc=2mm, boxrule=0.8pt,
  left=2mm, right=2mm, top=1.5mm, bottom=1.5mm,
  title={\footnotesize\bfseries\color{white}\faSitemap\ SƠ ĐỒ TƯ DUY: TIẾN TRÌNH 5 BƯỚC NHÂN PHẢI SANG TRÁI}
]
\centering
\begin{tikzpicture}[
  box/.style={draw, rounded corners=2pt, align=center, font=\scriptsize\bfseries, inner sep=2.5pt, minimum height=0.75cm},
  arr/.style={-{Stealth[length=2mm]}, thick}
]
  \node[box, fill=purple!15, draw=purple!80!black, text width=2.1cm] (n1) at (0, 0) {BƯỚC 1: ĐV\\$e \times k$\\Nhớ sang chục};
  \node[box, fill=teal!15, draw=teal!80!black, text width=2.1cm] (n2) at (2.8, 0) {BƯỚC 2: CHỤC\\$d \times k + \text{nhớ}$\\Nhớ sang trăm};
  \node[box, fill=blue!12, draw=blue!80!black, text width=2.1cm] (n3) at (5.6, 0) {BƯỚC 3: TRĂM\\$c \times k + \text{nhớ}$\\Nhớ sang nghìn};
  \node[box, fill=green!15, draw=green!70!black] (n4) at (8.4, 0) {BƯỚC 4: NGHÌN\\$b \times k + \text{nhớ}$\\Nhớ sang chục nghìn};
  \node[box, fill=orange!15, draw=orange!80!black, text width=2.3cm] (n5) at (11.2, 0) {CHỤC NGHÌN\\$a \times k + \text{nhớ}$\\Ghi kết quả};

  \draw[arr, purple!80!black] (n1) -- (n2);
  \draw[arr, teal!80!black] (n2) -- (n3);
  \draw[arr, blue!80!black] (n3) -- (n4);
  \draw[arr, green!70!black] (n4) -- (n5);
\end{tikzpicture}
\end{tcolorbox}

\vspace{-2.5mm}
% --- KHỐI 4: VÍ DỤ & MẸO TRÁNH BẪY ---
\noindent
\begin{minipage}[t]{0.49\textwidth}\vspace{0pt}
\begin{tcolorbox}[
  enhanced,
  colback=green!4!white,
  colframe=green!60!black,
  arc=2mm, boxrule=0.8pt,
  left=2mm, right=2mm, top=1.8mm, bottom=1.8mm,
  title={\footnotesize\bfseries\color{white}\faCheckCircle\ Ví dụ thực hành chuẩn mực}
]
{\footnotesize
\textbf{1. Đặt tính nhân có nhớ:}\\
$\begin{array}{r@{}l}
  16240 & \\
  \times\quad 4 & \\
  \hline
  \mathbf{64960} &
\end{array}$
\qquad
\begin{minipage}[c]{0.5\linewidth}
-- $4 \times 0 = 0$, viết $0$.\\
-- $4 \times 4 = 16$, viết $6$ nhớ $1$.\\
-- $4 \times 2 = 8$, thêm $1$ là $9$, viết $9$.\\
-- $4 \times 6 = 24$, viết $4$ nhớ $2$.\\
-- $4 \times 1 = 4$, thêm $2$ là $6$, viết $6$.
\end{minipage}

\textbf{2. Nhân ra số 100 000:}\\
$25\,000 \times 4 = \mathbf{100\,000}$.
}
\end{tcolorbox}
\end{minipage}
\hfill
\begin{minipage}[t]{0.49\textwidth}\vspace{0pt}
\begin{tcolorbox}[
  enhanced,
  colback=red!4!white,
  colframe=red!70!black,
  arc=2mm, boxrule=0.8pt,
  left=2mm, right=2mm, top=1.8mm, bottom=1.8mm,
  title={\footnotesize\bfseries\color{white}\faExclamationTriangle\ Mẹo tránh bẫy trắc nghiệm}
]
{\footnotesize
\begin{itemize}\setlength{\itemsep}{1.5pt}\setlength{\leftskip}{-2.5mm}
  \item \textbf{Bẫy nhân với chữ số 0 tận cùng:}\\
  $16\,240 \times 4$, nhiều bạn bỏ quên chữ số $0$ ở hàng đơn vị và ghi kết quả là $6496$! Kết quả đúng phải là $\mathbf{64\,960}$.
  \item Luôn nhớ: Nhân trước rồi mới cộng số nhớ vào sau!
\end{itemize}
}
\end{tcolorbox}
\end{minipage}


\end{document}
